\chapter{Grafi expander} \label{chap:expander}

% mi piacerebbe mettere una citazione all'inizio di ogni capitolo
% "the hay in the haystack"?

In questo capitolo introduciamo i grafi expander e le loro proprietà
pseudo-casuali. L'argomento è estremamente vasto, noi ne copriamo poco più di
quanto basta per dimostrare il Teorema PCP e per migliorare il risultato di
inapprossimabilità di \texttt{MAX-CLIQUE}.
Questo è l'unico capitolo in cui non presentiamo tutte le dimostrazioni dei
risultati che esulano da un tipico programma triennale di matematica o
informatica, per le quali si rimanda per esempio a \cite{Tre16, ABbook}.
L'unico fatto che nel Capitolo~\ref{chap:pcpproof} useremo come scatola nera è
l'esistenza di un algoritmo che per ogni $n$ costruisce un expander di $n$ nodi
in tempo polinomiale in $n$.

\section{Grafi e autovalori}

Un grafo si dice \emph{regolare di grado $d$}, d'ora in poi $d$-regolare, se
ogni nodo ha esattamente $d$ vicini. Per comodità, questo capitolo tratterà
esclusivamente grafi non diretti e $d$-regolari.

La matrice di adiacenza di un grafo $G$ di $n$ vertici è la matrice
$n \times n$ che in posizione $(u, v)$ ha il numero di archi di $G$
da $u$ a $v$.
Per un grafo non diretto e $d$-regolare $G$, indichiamo con $A = A(G)$ la sua
matrice di adiacenza rinormalizzata in modo tale che le entrate su ogni riga
(e colonna) sommino a $1$. Equivalentemente, $A$ è la matrice di adiacenza di
$G$ moltiplicata per lo scalare $\frac1{d}$, oppure la matrice di transizione
delle passeggiate aleatorie sul grafo.

Per il teorema spettrale, la matrice $A$, che è simmetrica, ammette una base
ortogonale di autovettori, i cui corrispondenti autovalori indichiamo con
$\lambda_1 \ge \lambda_2 \ge \dots \ge \lambda_n$.
Detto $\one = \big(\frac1{n}, \dots, \frac1{n}\big) \in \R^n$ il vettore costante
$\frac1{n}$ in tutte le entrate, vale $A \one = \one$, dunque $\one$ è
autovettore di $A$ relativo all'autovalore $1$.
Ciò corrisponde al fatto che la distribuzione uniforme sui nodi del grafo è
invariante per passeggiate aleatorie grazie all'ipotesi di $d$-regolarità.
Valgono inoltre\footnote{Per esempio per il primo teorema di Ger\v{s}gorin.}
$\lambda_1 = 1$ e $\lambda_n \ge -1$.

Saremo particolarmente interessati al modulo del secondo autovalore di modulo
massimo.

\begin{definition}
	Siano $G$ come sopra un grafo $d$-regolare non diretto e $A$ la sua matrice
	di adiacenza rinormalizzata.
	Definiamo il parametro
	$\lambda(G) := \max\{\labs \lambda_2 \rabs, \labs \lambda_n \rabs\}$.
\end{definition}

Poiché la matrice $A$ è simmetrica, il modulo del secondo autovalore massimo
in modulo ammette un'equivalente definizione variazionale.

\begin{proposition}
	Vale
	\[
		\lambda(G) =
		\max_{\substack{\xx \perp \one \\ \xx \neq 0}}
		\frac{\norm{A \xx}}{\norm{\xx}},
	\]
	dove la norma è quella euclidea.
\end{proposition}

Questa caratterizzazione sarà senza alcun dubbio la più utile.
Siamo pronti per introdurre finalmente gli expander. La definizione che segue
non è l'unica possibile, come vedremo nel corso del capitolo.

\begin{definition}[Grafi expander]
	Fissato $\lambda < 1$, chiamiamo \emph{$(n, d, \lambda)$-expander} un grafo
	di $n$ nodi, regolare di grado $d$ e con $\lambda(G) \le \lambda$.
	Definiamo allo stesso modo $(n, \lambda)$-expander e
	$(d, \lambda)$-expander nei casi in cui non specifichiamo il grado o il
	numero di vertici.
\end{definition}

\section{Un expander si comporta come un grafo casuale}

In questa sezione presentiamo alcune delle proprietà pseudo-casuali degli
expander. I seguenti risultati confermano la bontà della scelta di $\lambda(G)$
come misura del comportamento pseudo-casuale del grafo.
Il primo tra questi mostra che l'estremo di una passeggiata aleatoria su un
expander converge molto velocemente alla distribuzione uniforme.

\begin{theorem} \label{thm:expander-conv}
	Sia $G$ un $(n, \lambda)$-expander e $\pp$ una distribuzione di probabilità
	sui nodi di $G$. Allora
	\[
		\norm{A^\ell \pp - \one} \le \lambda^\ell.
	\]
\end{theorem}

In particolare, il Teorema \ref{thm:expander-conv} implica che l'estremo di una
passeggiata aleatoria di $\ell$ passi su un expander è distribuito quasi
uniformemente già per $\ell = O(\log n)$.
Avremmo potuto definire equivalentemente gli expander come i grafi per cui
vale la conclusione del Teorema~\ref{thm:expander-conv}\footnote{A priori non
è ovvio, poiché il vettore $\pp$ deve essere una distribuzione. Lasciamo la
dimostrazione come esercizio per il lettore.}.

\begin{proof}
	Vale $A^\ell \pp - \one = A^\ell (\pp - \one) = A^\ell \pp^\perp$, dove
	$\pp^\perp$ è la componente di $\pp$ ortogonale a $\one$. Sia ha dunque
	$\norm{A^\ell \pp^\perp} \le \lambda^\ell \norm{\pp^\perp}$.
	Per concludere è sufficiente maggiorare la norma di $\pp^\perp$ con quella
	di $\pp$ e applicare la disuguaglianza tra norma $2$ e norma $1$ in
	$\R^n$: $\norm{\pp^\perp} \le \norm{\pp} \le \lvert \pp \rvert_1 = 1$.
\end{proof}

Il seguente teorema mostra invece che è poco probabile i nodi di una
passeggiata aleatoria su un expander restino confinati in una piccola regione
del grafo.

\begin{theorem}[Passeggiate aleatorie su expander] \label{thm:expander-walk}
	Sia $G$ un $(n, d, \lambda)$-expander e sia $B \subseteq V$ un sottoinsieme
	dei vertici che soddisfa $|B| \le \beta n$ per $\beta \in (0, 1)$.
	Detti $X_1, \dots, X_k$ i nodi di una passeggiata aleatoria su $G$ con $X_1$
	scelto uniformemente in $V$, vale
	\[
		\Pr[\forall 1 \le i \le k \ X_i \in B]
		\le ((1 - \lambda)\sqrt{\beta} + \lambda)^{k - 1}.
	\]
\end{theorem}
\begin{proof}[Dimostrazione del Teorema~\ref{thm:expander-walk}]
	Esprimiamo algebricamente la probabilità alla sinistra dell'uguale.
	Con un abuso di notazione, indichiamo con $B$ la proiezione $\R^n \to \R^n$
	sulle coordinate in $B \subseteq V$. La probabilità sulla sinistra si
	scrive allora come $\lvert (B A)^{k - 1} B \one \rvert_1$, dove
	$(B A)^{k - 1} B \one$ è il vettore che, nell'entrata $j$, ha la probabilità
	che una passeggiata aleatoria lunga $k - 1$ termini in $j$ e non esca mai
	da $B$.
	Vale dunque
	\begin{align*}
		\Pr[\forall 1 \le i \le k \ X_i \in B] &=
			\lvert (B A)^{k - 1} B \one \rvert_1 \\
		&\le \sqrt{n} \norm{(B A)^{k - 1} B \one} \\
		&\le \sqrt{n} \norm{(B A)^{k - 1}} \norm{B \one} \\
		&= \sqrt{n} \norm{(B A)^{k - 1}} \sqrt{\beta} / \sqrt{n} \\
		&\le \norm{B A}^{k - 1}.
	\end{align*}
	L'obiettivo ora è maggiorare la norma matriciale di $B A$ con la quantità
	$(1 - \lambda)\sqrt{\beta} + \lambda$, per cui sarà	utile il seguente
	\begin{lemma}
		Sia $A$ come sopra. Allora
		\[
			A = (1 - \lambda) J + \lambda C,
		\]
		con $\norm{C} \le 1$ e $J_{i j} = \frac1{n}$ per ogni $i,j$.
	\end{lemma}
	\begin{proof}
		La matrice $C$ è interamente determinata da $A$, nello specifico vale
		$C = \frac1{\lambda}\big( A - (1 - \lambda)J \big)$.
		Per stimare la norma matriciale di $C$, come abbiamo già fatto,
		scomponiamo un generico vettore $\vv \in \R^n$ nelle componenti
		parallela e perpendicolare a $\one$.
		Vale $C \vv = C (\vv^\parallel + \vv^\perp)
		= \vv^\parallel + \frac1{\lambda} A \vv^\perp$.
		Ancora una volta applichiamo la definizione variazionale di $\lambda$
		e concludiamo.
	\end{proof}

	È sufficiente scrivere la matrice $A$ come nella tesi del lemma appena
	dimostrato, notare $\norm{B J} = \sqrt{\beta}$ e concludere per
	subadditività della norma:
	\begin{align*}
		\norm{B A} &= \big\lVert B \big( (1 - \lambda) J
			+ \lambda C \big) \big\rVert \\
		&\le (1 - \lambda) \norm{B J} + \lambda \norm C \\
		&\le (1 - \lambda) \sqrt{\beta} + \lambda.
	\end{align*}
\end{proof}

Il Teorema~\ref{thm:expander-walk} è alla base del ``riciclo di bit casuali'',
che vedremo applicato nella dimostrazione del
Corollario~\ref{thm:maxclique-inapprox-nc}.

I seguenti due lemmi, invece, saranno le espressioni di pseudocasualità degli
expander di cui faremo uso nel Capitolo~\ref{chap:pcpproof} per la
dimostrazione del Teorema~\ref{thm:pcp}.

\begin{restatable}{lemma}{expanderFwalk} \label{thm:Din06Prop2.5}
	% serve without self-loops? No, purché un cappio abbia contributo 2 al grado
	Sia $G = (V, E)$ un $(d, \lambda)$-expander e $F \subseteq E$ un sottoinsieme
	degli archi. Sia $K$ la distribuzione di probabilità sui vertici indotta
	dalla scelta di un estremo di un arco in $F$ scelto uniformemente a caso.
	La probabilità $p$ che una passeggiata aleatoria che comincia con
	distribuzione $K$ faccia l'$(i+1)$-esimo passo in $F$ è maggiorata da
	$\frac{|F|}{|E|} + \lambda^i$.
\end{restatable}

\begin{proof}
	Siano $\xx$ il vettore distribuzione di $K$, cioè $x_v = \frac{\deg_F(v)}
	{2 |F|}$, e $y_v$ la probabilità che un arco incidente a $v$ sia in $F$, cioè
	$\yy = \frac{2 |F|}{d} \xx$.
	Scomponiamo $\xx = \xx^\parallel + \xx^\perp$ con
	$\xx^\parallel := \frac1n \one$. Vale per definizione
	\[
		p = \sum_{v \in V}{ y_v (A^i \xx)_v }
		  = \langle \yy, A^i \xx \rangle.
	\]

	Osserviamo $\norm{\xx} ^2 \le \big(\sum_v{\lvert x_v \rvert} \big)
	\max_v{ \lvert x_v \lvert } = \max_v{ \lvert x_v \rvert }
	\le \frac{d}{2 | F |}$ da cui, applicando la disuguaglianza di Cauchy-Schwarz
	e la definizione variazionale di $\lambda$, otteniamo
	\[
		\big\langle \yy, A^i \xx^\perp \big\rangle
		\le \norm{\yy} \cdot \norm{A^i \xx^\perp}
		\le \frac{2 |F|}{d} \norm{\xx} \cdot \lambda^i \norm{x}
		\le \lambda^i.
	\]

	In conclusione,
	\[
		\langle \yy, A^i \xx \rangle
		= \big\langle \yy, A^i \xx^\parallel \big\rangle
		+ \big\langle \yy, A^i \xx^\perp \big\rangle
		\le \frac{2 \lvert F \rvert}{n d} + \lambda^i.
	\]
\end{proof}

\begin{lemma} \label{thm:ESbarS}
	Sia $G = (V, E)$ un $(n, d, \lambda)$-expander e $S \subseteq V$ un
	sottoinsieme dei vertici. Allora
	\[
		\lvert E(S, \bar{S}) \rvert \ge (1 - \lambda)\frac{d}{n}|S||\bar{S}|,
	\]
	dove con $E(S, \bar S)$ indichiamo l'insieme degli archi tra $S$ e $\bar S$.
\end{lemma}
\begin{proof}
	Esprimiamo $\lvert E(S, \bar S) \rvert$ come un opportuno prodotto scalare e
	analizziamo separatamente componente parallela e perpendicolare al vettore
	costante.
	Detta $A$ la matrice di adiacenza rinormalizzata, siano $\bm{1}_S$ e
	$\bm{1}_{\bar{S}}$ i vettori indicatori di $S$ e del suo complementare.
	Siano anche $\bm{1}_S^\parallel = \frac{|S|}{n} \bm{1}$ e
	$\bm{1}_S^\perp = \bm{1}_S - \bm{1}_S^\parallel$ le componenti
	di $\bm{1}_S$ parallela e perpendicolare a $\bm{1}$.
	Vale
	\begin{align*}
		\frac{1}{d} \lvert E(S, \bar S) \rvert
			&= \langle \bm{1}_S, A \bm{1}_{\bar S} \rangle \\
		&= \langle \bm{1}_S, A (\bm{1} - \bm{1}_{\bar S}) \rangle \\
		&= \langle \bm{1}_S, \bm{1} \rangle
			- \langle \bm{1}_S, A \bm{1}_S \rangle \\
		&= |S| - \big\langle \bm{1}_S^\parallel + \bm{1}_S^\perp,
			A (\bm{1}_S^\parallel + \bm{1}_S^\perp ) \big\rangle \\
		&= |S| - \big\langle \bm{1}_S^\parallel, \bm{1}_S^\parallel \big\rangle
			- \big\langle \bm{1}_S^\perp, A \bm{1}_S^\perp \big\rangle
			- 2 \big\langle \bm{1}_S^\parallel, \bm{1}_S^\perp \big\rangle \\
		&\ge |S| - \big\langle \bm{1}_S^\parallel, \bm{1}_S^\parallel \big\rangle
			- \lambda \big\langle \bm{1}_S^\perp, \bm{1}_S^\perp \big\rangle \\
		&= |S| - \frac{|S|^2}{n} - \lambda\bigg(|S| - \frac{|S|^2}{n}\bigg) \\
		&= (1 - \lambda)\frac{|S|}{n}(n - |S|) = (1 - \lambda) \frac1{n} |S| |\bar S|.
	\end{align*}
\end{proof}

Osserviamo che la quantità $\frac{d}{n}|S| |\bar S|$ corrisponde essenzialmente
al valore atteso del numero di archi tra $S$ e $\bar S$ in un grafo
$d$-regolare scelto uniformemente a caso.
Ignorando per semplicità i cappi, su $\binom{n}{2}$ possibili archi esattamente
$|S| |\bar S|$ contengono un estremo in $S$ e l'altro in $\bar S$.
Poiché un grafo $d$-regolare ha $\frac{n d}{2}$ archi, per simmetria ognuno di
questi appare in una frazione
$\frac{n d}{2} {\binom{n}{2}}^{-1} = \frac{d}{n - 1}$ dei grafi $d$-regolari su
$n$ nodi, da cui $\E[\lvert E(S, \bar S) \rvert]
\approx \frac{d}{n} |S| |\bar S|$.
Il Lemma~\ref{thm:ESbarS} garantisce che il numero di archi tra $S$ ed $\bar S$
in un expander $G$ non è molto minore del numero di archi tra $S$ ed $\bar S$
in un grafo casuale.

\section{Un grafo casuale è con alta probabilità un expander}

Fin'ora abbiamo osservato che, se il parametro $\lambda(G)$ di un grafo
$d$-regolare è distante da $1$ almeno una costante indipendente dalla
dimensione del grafo, allora questo somiglia a un grafo casuale.
In questa sezione ci occupiamo del viceversa: mostriamo in particolare che un
grafo generato tramite un opportuno processo probabilistico è, con alta
probabilità, un expander.
Ciò dimostrerà l'esistenza di expander, che a priori non scontata.

Per percorrere questa strada è conveniente dare una terza definizione
equivalente di expander come, a meno di cambiare nome ai parametri, i grafi per
cui vale il Lemma~\ref{thm:ESbarS}.

\begin{definition}
	Fissato $\rho > 0$, chiamiamo \emph{$(n, d, \rho)$-expander combinatorio}
	un grafo di $n$ nodi, regolare di grado $d$ e tale per cui, per ogni
	sottoinsieme $S$ di al più $n / 2$ nodi, vale
	\[
		\lvert E(S, \bar S) \rvert \ge \rho d \lvert S \rvert.
	\]
\end{definition}

Il Lemma~\ref{thm:ESbarS} garantisce che un $(n, d, \lambda)$-expander è un
$(n, d, \rho)$-expander combinatorio per $\rho = \frac{1 - \lambda}{2}$.
Vale anche che un expander combinatorio è un expander, ma la dimostrazione
richiede un po' più di lavoro. Rimandiamo per esempio al Capitolo 21 di
\cite{ABbook}.

Siamo pronti per descrivere un processo casuale che, con alta probabilità,
produce un expander combinatorio, quindi un expander.

Supponiamo per semplicità $n$ pari e, per $d \in \N$, consideriamo un grafo
$d$-regolare generato dal seguente processo.
Fissato l'insieme di vertici $V$ con $|V| = n$, siano $M_1, \dots, M_d$
matching perfetti scelti uniformemente a caso.
Definiamo il grafo casuale $G = (V, E)$ come l'unione dei $d$ matching:
\[
	E = M_1 \cup \dots \cup M_d.
\]
Il grafo $G$ è regolare di grado $d$ e privo di cappi.

È utile pensare a un matching distribuito uniformemente come generato tramite
il seguente processo\footnote{La costruzione e la dimostrazione che seguono
ricalcano le lezioni di Trevisan \cite{Tre16}.}.
Fissiamo un ordinamento arbitrario dei vertici. A ogni passo, consideriamo
il primo vertice $v$ non ancora accoppiato e scegliamo il suo compagno $w$
uniformemente a caso tra i vertici rimanenti. Il processo termina quando
tutti i vertici sono stati accoppiati.
È facile verificare che questo algoritmo produce un matching distribuito
uniformemente tra gli $(n - 1) \cdot (n - 3) \cdots \cdot 1$ matching perfetti
su $V$.


\begin{theorem} \label{thm:random-expander}
	Siano $n$ un intero pari e $G$ il grafo su $n$ vertici ottenuto dall'unione
	di $d = 25$ matching perfetti, scelti in modo indipendente e con
	distribuzione uniforme.
	Allora, con probabilità $1 - O\big(\frac{1}{n^4}\big)$, il grafo $G$ è un
	$\big(n, 25, \frac{1}{150}\big)$-expander combinatorio.
\end{theorem}

\begin{proof}
	Mostriamo che, con alta probabilità, ogni sottoinsieme $S \subseteq V$ con
	$|S| \le n / 2$ soddisfa $|E(S, \bar S)| \ge \rho d |S|$ per $d = 25$ e
	$\rho = \frac{1}{150}$. Questo equivale a richiedere
	\[
		|E(S, \bar S)| \ge \frac{|S|}{6}.
	\]
	Indichiamo con $\Gamma(S)$ l'insieme dei vertici adiacenti ad almeno un
	nodo di $S$ in $G$. Poiché ogni nodo in $\Gamma(S) \setminus S$ è
	raggiunto da almeno un arco con un estremo in $S$, si ha
	$|E(S, \bar S)| \ge |\Gamma(S) \setminus S|$.
	Di conseguenza, se un insieme $S$ non soddisfa la proprietà di espansione,
	allora $|\Gamma(S) \setminus S| < |S| / 6$.
	In tal caso, deve esistere un sottoinsieme $T \subseteq \bar S$ di cardinalità
	$|T| \le |S|/6$ tale che $\Gamma(S) \setminus S \subseteq T$, ovvero tale per cui
	$\Gamma(S) \subseteq S \cup T$.
	Mostreremo che, per ogni $T$ fissato, ciò è molto improbabile, poi
	concluderemo con due limiti per l'unione.
	% TBD. Meglio union bound?
	
	Fissiamo allora un sottoinsieme $S$ di cardinalità $|S| = k \le n/2$ e un
	sottoinsieme $T \subseteq \bar S$ di cardinalità al più $k/6$ (possiamo supporre,
	senza perdita di generalità, $|T| = \lfloor k/6 \rfloor$).
	Maggioriamo la probabilità che in un singolo matching perfetto $M$ si verifichi
	$\Gamma_M(S) \subseteq S \cup T$.
	Possiamo supporre il matching $M$ sia generato dal processo iterativo
	descritto sopra fissato un'ordinamento dei vertici in cui gli elementi di
	$S$ precedono tutti gli altri.
	Finché ci sono nodi di $S$ da accoppiare, il vertice $v$ estratto
	per trovare un compagno appartiene a $S$. Poiché in una singola iterazione
	al massimo due nodi di $S$ vengono accoppiati, sono necessarie almeno $k/2$
	iterazioni per accoppiare l'intero insieme $S$.
	La condizione $\Gamma_M(S) \subseteq S \cup T$ implica che, in ciascuna di
	queste prime iterazioni, il compagno $w$ selezionato per un nodo di $S$ deve
	necessariamente appartenere all'insieme $S \cup T$.
	
{\color{magenta}
	Consideriamo l'$i$-esima iterazione ($1 \le i \le k/2$) e condizioniamo al
	fatto che, nei passi precedenti, siano stati scelti esclusivamente nodi in
	$S \cup T$. In questo momento sono già stati accoppiati $2(i - 1)$ vertici,
	tutti appartenenti a $S \cup T$.
	Di conseguenza, i nodi di $S \cup T$ ancora disponibili come potenziali
	partner sono al più
	\[
		|S \cup T| - (2i - 1) \le \frac{7}{6}k - 2i + 1.
	\]
	D'altro canto, il numero totale di vertici liberi nel grafo è $n - 2i + 1$.
	Pertanto, la probabilità che all'$i$-esimo passo si scelga un partner $w \in S \cup T$
	è limitata da
	\[
		\frac{\frac{7}{6}k - 2i + 1}{n - 2i + 1}.
	\]
	Il verificarsi dell'evento $\Gamma_M(S) \subseteq S \cup T$ richiede che questa
	condizione valga per ciascuna delle prime $k/2$ iterazioni. Moltiplicando le
	rispettive probabilità condizionate, otteniamo
	\[
		\Pr[\Gamma_M(S) \subseteq S \cup T] \le \prod_{i=1}^{k/2} \frac{\frac{7}{6}k - 2i + 1}{n - 2i + 1}.
	\]
	Per maggiorare questo prodotto in modo maneggevole, dividiamo i fattori in due gruppi.
	Per $i > k/3$, il numeratore soddisfa
	\[
		\frac{7}{6}k - 2i + 1 \le \frac{7}{6}k - \frac{2}{3}k = \frac{1}{2}k.
	\]
	Al contempo, essendo $k \le n/2$ e $i \le k/2$, il denominatore è limitato inferiormente da
	\[
		n - 2i + 1 > n - k \ge n - \frac{n}{2} = \frac{1}{2}n.
	\]
	Ne consegue che, per ciascuno degli ultimi $k/2 - k/3 = k/6$ fattori, il rapporto
	è maggiorato da $\frac{k/2}{n/2} = \frac{k}{n}$. Maggiorando per semplicità con $1$ i
	primi $k/3$ fattori della produttoria, otteniamo
	\[
		\Pr[\Gamma_M(S) \subseteq S \cup T] \le \left(\frac{k}{n}\right)^{k/6}.
	\]
	Passando ora all'intero grafo $G$, poiché i $d$ matching $M_1, \dots, M_d$ sono estratti
	in modo indipendente, l'evento $\Gamma(S) \subseteq S \cup T$ si realizza solo se esso si
	verifica in tutti i $d$ matching. Pertanto,
	\[
		\Pr[\Gamma(S) \subseteq S \cup T] = \big(\Pr[\Gamma_M(S) \subseteq S \cup T]\big)^d \le \left(\frac{k}{n}\right)^{\frac{dk}{6}}.
	\]
	Dobbiamo ora fare un \emph{union bound} su tutte le possibili scelte per gli insiemi
	$S$ e $T$. Ricordando la disuguaglianza $\binom{n}{m} \le \big(\frac{en}{m}\big)^m$, valida
	per ogni $1 \le m \le n$, limitiamo il numero di scelte:
	\[
		\binom{n}{k} \le \left(\frac{en}{k}\right)^k, \qquad
		\binom{n}{\lfloor k/6 \rfloor} \le \left(\frac{6en}{k}\right)^{k/6}.
	\]
	Così, il numero di possibili coppie $(S, T)$ per un fissato $k$ è
	maggiorato da
	\[
		\binom{n}{k} \binom{n}{\lfloor k/6 \rfloor}
		\le e^k (6e)^{k/6} \left(\frac{n}{k}\right)^{\frac{7k}{6}}.
	\]
	Definiamo la costante $C := e (6e)^{1/6} \approx 4.33$. La probabilità che
	esista \emph{almeno} una coppia $(S, T)$ della dimensione considerata per cui fallisce
	l'espansione è
	\begin{align*}
		\Pr\big[\exists S, T : |S|=k, |T| \le k/6, \Gamma(S) \subseteq S \cup T\big]
		&\le e^k (6e)^{k/6} \left(\frac{n}{k}\right)^{\frac{7k}{6}} \left(\frac{k}{n}\right)^{\frac{dk}{6}} \\
		&= \left[ C \left(\frac{k}{n}\right)^{\frac{d - 7}{6}} \right]^k.
	\end{align*}
	Sostituendo $d = 25$, l'esponente all'interno della parentesi diventa $\frac{25 - 7}{6} = 3$.
	Poiché abbiamo assunto $k \le n/2$, segue che $\frac{k}{n} \le \frac{1}{2}$. Questo garantisce che
	la base della potenza sia strettamente minore di $1$:
	\[
		C \left(\frac{k}{n}\right)^3 \le 4.33 \cdot \left(\frac{1}{2}\right)^3 = \frac{4.33}{8} < 0.55 < 1.
	\]
	Rimane da sommare questa probabilità di fallimento su tutte le possibili dimensioni di $S$,
	ossia per $1 \le k \le n/2$:
	\begin{itemize}
		\item Per $k = 1$, il grafo è per costruzione regolare di grado $d = 25$ e privo di
		cappi. L'unico vertice $v \in S$ deve quindi avere $25$ archi incidenti verso $\bar S$,
		perciò $|E(S, \bar S)| = 25 > \frac{1}{6} \cdot 1$. La condizione di espansione per insiemi
		di taglia $1$ è dunque soddisfatta con certezza.
		\item Per $k = 2$, la probabilità di fallimento è limitata da
		\[
			\left[ C \left(\frac{2}{n}\right)^3 \right]^2 = \frac{64 C^2}{n^6} = O\left(\frac{1}{n^6}\right).
		\]
		\item Per $3 \le k \le \lfloor\sqrt{n}\rfloor$, sfruttando l'ipotesi $\frac{k}{n} \le \frac{1}{\sqrt{n}}$,
		la base risulta maggiorata da $C(k/n)^3 \le C n^{-3/2}$, da cui
		\[
			\sum_{k=3}^{\lfloor\sqrt{n}\rfloor} \left[ C \left(\frac{k}{n}\right)^3 \right]^k
			\le \sum_{k=3}^{\lfloor\sqrt{n}\rfloor} \left(\frac{C}{n^{3/2}}\right)^k
			\le \sum_{k=3}^{\lfloor\sqrt{n}\rfloor} \frac{C^k}{n^{3k/2}}
			= O\left(\frac{1}{n^{9/2}}\right).
		\]
		\item Per $\lfloor\sqrt{n}\rfloor < k \le n/2$, la base è uniformemente
		maggiorata dalla costante $0.55$, per cui
		\[
			\sum_{k = \lfloor\sqrt{n}\rfloor + 1}^{n/2} \left[ C \left(\frac{k}{n}\right)^3 \right]^k
			\le \sum_{k > \sqrt{n}} (0.55)^k
			= O\big(0.55^{\sqrt{n}}\big) = o\left(\frac{1}{n^c}\right) \quad \text{per ogni } c > 0.
		\]
	\end{itemize}
	Combinando queste disuguaglianze, la probabilità complessiva che esista un insieme che viola
	l'espansione è dominata asintoticamente dal termine per $k=2$, e vale pertanto
	$O(1/n^6) + O(1/n^{4.5}) + o(1/n^c) = O(1/n^4)$.
	
	Concludiamo che, con probabilità perlomeno $1 - O(1/n^4)$, ogni insieme $S$ con
	$1 \le |S| \le n/2$ soddisfa il bound $|E(S, \bar S)| \ge \frac{1}{6}|S| = \frac{1}{150} \cdot 25 |S|$.
	Questo prova che $G$ è un $\big(n, 25, \frac{1}{150}\big)$-expander combinatorio.
}
\end{proof}

\section{Costruzioni esplicite}

\begin{definition}[Famiglie di expander]
	Una famiglia di grafi $\{G_n\}_{n \in \N}$ si dice \emph{famiglia di
	$(d, \lambda)$-expander} se esistono costanti $d \in \N$ e $\lambda < 1$
	tali che, per ogni $n$, il grafo $G_n$ è un $(n, d, \lambda)$-expander.
	
	Chiameremo famiglia di expander anche una famiglia di grafi $G_n$ come sopra
	definiti solo per $n \ge n_0$.
\end{definition}

\lorenzo{Lascio Zig-Zag? Forse no, meglio citarlo e basta.}

\begin{theorem} \label{thm:expander-construction}
	Esiste una famiglia fortemente esplicita di $(d, \lambda)$-expander per
	qualche $d \in \mathbb{N}$ e $\lambda < 1$.
\end{theorem}
\begin{proof}
	\lorenzo{Al massimo a grandi linee.}
\end{proof}

\lorenzo{Voglio definire l'elevamento a potenza?}

Notiamo che grazie all'elevamento a potenza di grafi il Teorema
\ref{thm:expander-construction} fornisce famiglie fortemente esplicite di
$\lambda$-expander per ogni $\lambda > 0$ al prezzo di un aumento del grado $d$.

\begin{corollary} \label{thm:maxclique-inapprox-nc}
	Esiste una costante $c > 0$ tale che non esiste una
	$\frac1{n^c}$-approssimazione polinomiale di \texttt{MAX-CLIQUE}, a meno
	che $\P = \NP$.
\end{corollary}

% citare https://lucatrevisan.github.io/pcp/lecture05.pdf

\begin{proof}
	Stessa riduzione del Teorema~\ref{thm:max-clique}, ma stavolta il
	verificatore PCP si ripete $\log n$ volte.
\end{proof}

